\chapter{Gait Disorder (Abnormal Gait)}

\section*{Definition}
Gait disorder refers to abnormal walking patterns due to neurologic, musculoskeletal, or systemic pathology. It is broadly classified into hemiparetic, ataxic, parkinsonian, neuropathic (steppage), myopathic (waddling), and spastic gait patterns. Gait impairment is a major risk factor for falls and loss of independence in older adults.

\section*{Pathophysiology}
Normal gait requires integration of the motor cortex, basal ganglia, cerebellum, brainstem locomotor centers, spinal central pattern generators, and peripheral sensory feedback. Cerebellar dysfunction (ataxia) disrupts coordination and truncal stability. Basal ganglia pathology (parkinsonism) produces bradykinesia and impaired postural reflexes. Lumbar stenosis or peripheral neuropathy reduces proprioceptive input, causing cautious or sensory ataxic gait. Upper motor neuron lesions produce spasticity and circumduction due to loss of descending inhibition. Muscle weakness shifts the center of mass over the stance limb, producing waddling or Trendelenburg gait.

\section*{Causes}
\begin{itemize}
  \item \textbf{Hemiparetic:} Stroke, brain tumor, multiple sclerosis, traumatic brain injury
  \item \textbf{Parkinsonian:} Parkinson disease, atypical parkinsonism (MSA, PSP, vascular parkinsonism)
  \item \textbf{Ataxic:} Cerebellar degeneration (alcohol, paraneoplastic, spinocerebellar ataxia), stroke, MS
  \item \textbf{Neuropathic (steppage):} Common fibular neuropathy, L5 radiculopathy, peripheral neuropathy, Charcot-Marie-Tooth
  \item \textbf{Myopathic (waddling):} Muscular dystrophy, polymyositis, inclusion body myositis, glucocorticoid myopathy
  \item \textbf{Splastic:} Cerebral palsy, spinal cord injury, multiple sclerosis, hereditary spastic paraparesis
  \item \textbf{Sensory:} Peripheral neuropathy, vitamin B12 deficiency, posterior column loss (tabes dorsalis)
  \item \textbf{Antalgic:} Hip or knee osteoarthritis, fracture, gout
  \item \textbf{Other:} Normal pressure hydrocephalus (magnetic gait), frontal lobe gait disorder, vestibular gait, psychogenic gait
\end{itemize}

\section*{When to See a Doctor}
See your doctor if:
\begin{itemize}
  \item New-onset gait difficulty, especially in older adults
  \item Gait impairment with recent falls or fear of falling
  \item Acute or subacute gait change with focal weakness, sensory loss, or back pain
  \item Gait abnormality with urinary incontinence or cognitive decline
\end{itemize}

\section*{Related Symptoms}
\begin{itemize}
  \item Circumduction (leg swings outward in a semicircle)
  \item Shuffling, festinating, or freezing gait
  \item Widened base of support and staggering (ataxia)
  \item Foot slapping or high-stepping (foot drop)
  \item Waddling (pelvic drop with each step)
  \item Difficulty with gait initiation or turning
\end{itemize}
\textbf{Important:} Normal pressure hydrocephalus should be suspected in the classic triad of gait apraxia (magnetic gait), urinary incontinence, and cognitive decline, especially if symptoms are of recent onset.

\section*{Self-Care / Home Management}
\begin{itemize}
  \item Use of appropriate walking aid (cane, walker, rollator) for stability
  \item Fall-proofing the home: remove rugs, improve lighting, install grab bars
  \item Wear sturdy, non-slip footwear
  \item Avoid multitasking while walking
  \item Daily walking within safe limits to maintain mobility
\end{itemize}

\section*{How Your Doctor Will Evaluate You}
\begin{itemize}
  \item Gait observation and classification: watch patient walk (initiation, speed, stride length, base width, arm swing, turning)
  \item Neurological examination assessing strength, sensation, coordination, reflexes, and postural stability
  \item Timed Get Up and Go test (TUG) for objective functional assessment
  \item MRI brain and/or cervical spine for UMN or cerebellar signs
  \item MRI lumbar spine for cauda equina or spinal stenosis
  \item Laboratory evaluation: vitamin B12, TSH, HbA1c, vitamin D, RPR, ESR
\end{itemize}

\section*{Treatment Options}
\begin{itemize}
  \item Physical therapy with gait training, balance exercises, and assistive device fitting
  \item Treatment of underlying condition: levodopa for Parkinson disease, shunt for NPH, decompression for spinal stenosis
  \item Strengthening of hip abductors, knee extensors, and ankle dorsiflexors for myopathic gait
  \item Ankle-foot orthosis for foot drop
  \item Vestibular rehabilitation for vestibular gait disorders
  \item Occupational therapy for home safety evaluation and fall prevention
  \item Medication review for drugs contributing to gait impairment (sedatives,antipsychotics)
\end{itemize}

\clearpage
